\documentclass[tikz,border=1cm]{standalone}
\usepackage{xcolor}
\usepackage{amsmath,amssymb}

\usepackage[utf8]{inputenc}
\usepackage[spanish]{babel}
\usepackage[
    letterpaper,
    top=1.5cm,
    bottom=2cm,
    left=2cm,
    right=2cm
]{geometry}

\usepackage{graphicx}
\usepackage{fancyhdr}
\usepackage[table,xcdraw]{xcolor}
\usepackage{tikz}
\usetikzlibrary{shapes.geometric, arrows, positioning, fit, backgrounds, calc, babel}
\usepackage{ifpdf}
\usepackage{blindtext}
\usepackage[cmex10]{amsmath}
\usepackage{amssymb}
\usepackage{amsfonts}
\usepackage{float}
\usepackage{siunitx}
\usepackage{listings}
\usepackage{hyperref}
\usepackage{subcaption}
\usepackage{booktabs}
\usepackage{url}
\usepackage{wrapfig}
\usepackage{titlesec}
\usepackage[export]{adjustbox}
\usepackage{imakeidx}
\usepackage{parskip}
\usepackage[numbers]{natbib}

\usepackage{ragged2e}
\usepackage{setspace}
\usepackage{array}
\usepackage{longtable}

% Definición de un nuevo tipo de columna para evitar "Underfull \hbox" en tablas
\newcolumntype{L}[1]{>{\RaggedRight\arraybackslash}p{#1}}

\hypersetup{colorlinks=false,bookmarksopen=true,linkbordercolor={1 1 1}}

\newcommand{\rfigura}[1]{Fig.\ref{#1}}
\newcommand{\rtabla}[1]{Tabla \ref{#1}}

% Datos del informe
\newcommand{\Curso}{IEE2913 --- Diseño Eléctrico (Capstone)}
\newcommand{\TituloInforme}{Informe de Avance 1 (5\%)}
\newcommand{\NumeroGrupo}{10}
\newcommand{\NombreProyecto}{SupaClock: Wearable Biométrico Modular}
\newcommand{\Integrantes}{Tomás Avendaño, Benjamín Sepúlveda, Pablo Uribe}
\newcommand{\FechaEntrega}{7 de abril de 2026}

\lstset{
  basicstyle=\ttfamily\small,
  breaklines=true,
  breakatwhitespace=true,
  postbreak=\mbox{\textcolor{red}{$\hookrightarrow$}\space},
  captionpos=b,
  frame=single,
  numbers=left,
  numberstyle=\tiny\color{gray},
  language=C,
  keywordstyle=\color{blue}\bfseries,
  commentstyle=\color{gray}\itshape,
  stringstyle=\color{red!70!black}
}


\begin{document}
\begin{tikzpicture}[
                node distance=1.2cm and 1.5cm,
                sensor/.style={rectangle, rounded corners, draw=black, top color=white, bottom color=gray!20, text width=2.5cm, align=center, minimum height=1cm, font=\footnotesize},
                task/.style={rectangle, rounded corners, draw=black, thick, top color=white, bottom color=blue!10, text width=4cm, align=center, minimum height=1.5cm, font=\small},
                queue/.style={rectangle, draw=black, fill=yellow!20, text width=2.5cm, align=center, minimum height=0.8cm, font=\footnotesize\ttfamily},
                mutex/.style={circle, draw=black, fill=red!20, inner sep=2pt, font=\footnotesize\bfseries},
                arrow/.style={thick, ->, >=stealth},
                dashed_arrow/.style={thick, dashed, ->, >=stealth},
                legend/.style={rectangle, draw=black!50, fill=white, inner sep=5pt, font=\scriptsize}
            ]%
            % Nodos de Hardware
            \node[sensor] (adc) {ADC / DMA\\(ECG 100Hz)};
            \node[sensor, below=of adc] (i2c) {Bus I2C\\(IMU, PPG, Temp)};
            \node[sensor, below=of i2c] (btns) {Botones físicos\\(Interrupciones)};
            %
            % Mutex
            \node[mutex, right=0.5cm of i2c] (mut) {M};
            %
            % Nodos de Tareas FreeRTOS
            \node[task, right=1.5cm of adc] (t_acq) {\textbf{Tarea Adquisición}\\ \textit{Prioridad: Alta}\\(Muestreo de sensores)};
            \node[task, below=1cm of t_acq] (t_proc) {\textbf{Tarea Procesamiento}\\ \textit{Prioridad: Media}\\(TinyML, Pasos)};
            \node[task, below=1cm of t_proc] (t_gui) {\textbf{Tarea Interfaz Gráfica}\\ \textit{Prioridad: Media-Baja}\\(Actualización SPI)};
            \node[task, right=1.5cm of t_proc] (t_ble) {\textbf{Tarea Com. BLE}\\ \textit{Prioridad: Baja}\\(Empaquetado y envío)};
            %
            % Colas (Queues)
            \node[queue, right=0.3cm of t_acq, yshift=-1cm] (q_data) {Queue: Datos};
            \node[queue, right=0.3cm of t_proc, yshift=-1cm] (q_ble) {Queue: Tx BLE};
            %
            % Background / Boundary para ESP32
            \begin{scope}[on background layer]
                \node[rectangle, draw=black!50, dashed, inner sep=0.5cm, fit=(adc) (i2c) (btns) (t_acq) (t_proc) (t_gui) (t_ble) (q_data) (q_ble) (mut), fill=black!2] (esp32) {};
                \node[anchor=north west, font=\bfseries] at (esp32.north west) {Microcontrolador ESP32 (FreeRTOS)};
            \end{scope}
            %
            % Conexiones Hardware -> Tareas
            \draw[arrow] (adc) -- (t_acq);
            \draw[arrow] (i2c) -- (mut) -- (t_acq);
            \draw[dashed_arrow] (btns) -- node[above, font=\scriptsize] {ISR} (t_gui);
            %
            % Conexiones entre Tareas (IPC)
            \draw[arrow] (t_acq) |- (q_data);
            \draw[arrow] (q_data) -| (t_proc);
            \draw[arrow] (t_proc) |- (q_ble);
            \draw[arrow] (q_ble) -| (t_ble);
            %
            % Conexiones de visualización
            \draw[arrow] (t_proc) -- (t_gui);
            %
            % Salidas externas
            \node[sensor, right=1.5cm of t_ble] (phone) {Smartphone\\(App Gateway)};
            \node[sensor, right=1.5cm of t_gui] (lcd) {Pantalla LCD\\(ST7789)};
            %
            \draw[arrow, blue, thick] (t_ble) -- node[above, font=\scriptsize] {BLE} (phone);
            \draw[arrow] (t_gui) -- node[above, font=\scriptsize] {SPI} (lcd);
            %
            % === LEYENDA ===
            \node[legend, below=0.5cm of esp32] (leg) {
                \begin{tabular}{cl}
                    \tikz\draw[fill=gray!20] (0,0) rectangle (0.3,0.2);   & Periférico de hardware       \\
                    \tikz\draw[fill=blue!10] (0,0) rectangle (0.3,0.2);   & Tarea FreeRTOS               \\
                    \tikz\draw[fill=yellow!20] (0,0) rectangle (0.3,0.2); & Cola de mensajes (Queue)     \\
                    \tikz\draw[fill=red!20] (0,0) circle (0.1);           & Mutex (exclusión mutua)      \\
                    \tikz\draw[thick, ->] (0,0.1) -- (0.4,0.1);           & Flujo de datos               \\
                    \tikz\draw[thick, dashed, ->] (0,0.1) -- (0.4,0.1);   & Interrupción asíncrona (ISR) \\
                    \tikz\draw[blue, thick, ->] (0,0.1) -- (0.4,0.1);     & Enlace BLE inalámbrico       \\
                \end{tabular}
            };
        \end{tikzpicture}
\end{document}
